\structelem{ТЕРМИНЫ И ОПРЕДЕЛЕНИЯ}

В настоящем отчете о ВКР применяют следующие термины с
соответствующими определениями.

\vspace{6pt}

% Перечень терминологических статей по ГОСТ 7.32-2017, п. 6.14:
% слева без абзацного отступа --- термин, справа через тире --- его
% определение, без знаков препинания в конце строки.
\begin{description}[
    font=\normalfont\bfseries,
    leftmargin=0pt, labelindent=0pt, labelsep=0.4em,
    style=sameline, itemsep=4pt, topsep=0pt
]

\item[Алерт] {\normalfont ---
сформированное системой обнаружения событие безопасности,
сопровождаемое уровнем критичности, скорингом, типом
предположительной атаки и контекстом наблюдения, направляемое в
SOC для последующей обработки аналитиком}

\item[Аномальная сетевая активность] {\normalfont ---
поведение сетевых соединений, статистически отклоняющееся от
характерного для контролируемой инфраструктуры профиля нормы,
проявляющееся в признаках длительности, направления, объема и
последовательности соединений}

\item[Аномальный детектор] {\normalfont ---
программный модуль, формирующий непрерывный скоринг
\(s(x)\colon X \to \mathbb{R}\) и принимающий решение об
аномальности наблюдения \(x\) при превышении порога \(\tau\)}

\item[Дрейф распределений] {\normalfont ---
изменение совместного распределения \(p(x, y)\) или его компонент
во времени в продуктивной эксплуатации; источник деградации
аномального детектора}

\item[Калибровка вероятностей] {\normalfont ---
процедура согласования выходной вероятности классификатора
с эмпирической частотой положительного класса; в работе
используется температурное масштабирование}

\item[Конечно-автоматный агент] {\normalfont ---
программный компонент активного тестирования, формирующий поток
сетевых событий с управляемым переключением режимов норма/атака/
дрейф на основе вероятностной матрицы переходов конечного автомата
и библиотеки параметризованных шаблонов трафика}

\item[Несанкционированные действия в сети] {\normalfont ---
действия, нарушающие установленную политику безопасности
корпоративной информационной системы, в том числе попытки получения
неавторизованного доступа, эксфильтрации данных, исчерпания ресурсов,
скрытого латерального перемещения и развертывания вредоносного кода}

\item[Сетевой поток (flow)] {\normalfont ---
агрегированная единица сетевого трафика, описываемая совокупностью
характеристик (src/dst IP, src/dst port, protocol) в пределах
заданного временного интервала и характеризуемая длительностью,
числом переданных пакетов и байтов, межпакетными интервалами и
сопутствующими статистиками}

\item[Скользящее окно (sliding window)] {\normalfont ---
последовательность из \(W\) подряд идущих flow-записей,
используемая в качестве входа модели; параметры окна --- размер
\(W\) и шаг \(S\)}

\item[Шаблон атаки] {\normalfont ---
параметризованное описание характеристик потока сетевых событий,
соответствующих определенному классу несанкционированной активности
(DDoS, port scan, brute force, exploit payload, data exfiltration,
internal reconnaissance, worm lateral movement)}

\item[NIDS (Network-based Intrusion Detection System)]
{\normalfont --- сетевая система обнаружения вторжений,
осуществляющая анализ проходящего трафика для выявления попыток
несанкционированного доступа}

\item[NTA (Network Traffic Analysis)] {\normalfont ---
подкласс систем мониторинга сетевой безопасности, ориентированный
на анализ метаданных сетевых соединений flow-уровня
с применением статистических и машинно-обучаемых моделей}

\item[SOC (Security Operations Center)] {\normalfont ---
функциональная организационно-техническая структура мониторинга
и реагирования на инциденты информационной безопасности
корпоративного класса}

\end{description}
